\section{Gaussian field setting and exact joint error ledger}\label{sec:setting}

Evaluating a continuum volatility functional after discretising the monitoring window, projecting the
Gaussian driver onto finitely many coordinates, and truncating its Wiener chaos expansion couples three
sources of approximation error. This section fixes the Hilbert-space representation contract and states
the exact joint mean-square identity, its two-sided residual-covariance bounds, and the measure-axis
convergence rate that underpin the financial certificates developed in Section~\ref{sec:certificates}.

\subsection{Subordinated Gaussian functionals and the representation contract}\label{sec:contract}
Let $\cH$ be a real separable Hilbert space and let $W=(W(h))_{h\in\cH}$ be an isonormal Gaussian process.
A continuous map $s\mapsto e_s\in\cH$ on a compact interval $\mathcal{I}\subset\R$ defines the centred
Gaussian field $Y_s=W(e_s)$ with covariance $R(s,u)=\ip{e_s}{e_u}$ and variance $\sigma_s^2=R(s,s)$.
Fractional Brownian motion for $H\in(0,1)$, the Riemann--Liouville rough-volatility driver
$W_t^H=\sqrt{2H}\int_0^t(t-s)^{H-1/2}\dd W_s$, and the forward-variance Gaussian field
$Y_u=\sqrt{2H}\int_0^T(u-s)^{H-1/2}\dd W_s$ on $u\in[T,T+\Delta]$ all fall into this setting; only the
covariance $R$ enters the error identities below.

\begin{definition}[Subordinated functional]\label{def:output}
Let $f\colon\mathcal{I}\times\R\to\R$ be measurable with $f(s,\cdot)\in L^2(N(0,\sigma_s^2))$ and
$\sup_{s\in\mathcal{I}}\norm{f(s,Y_s)}_2<\infty$. For any finite signed Borel measure $\rho$ on
$\mathcal{I}$, the subordinated output is the $L^2$-valued Bochner integral
\[
  \Lambda_\rho=\int_{\mathcal{I}} f(s,Y_s)\,\rho(\dd s).
\]
\end{definition}

With $\He_m$ the probabilists' Hermite polynomials and $h_m=\He_m/\sqrt{m!}$, the normalised Hermite
coefficients $\hat c_m(s)=\E[f(s,\sigma_sZ)h_m(Z)]$ ($Z\sim N(0,1)$) give the Wiener chaos expansion
$X_s=f(s,Y_s)=\sum_{m\ge0}\hat c_m(s)h_m(Y_s/\sigma_s)$. For $|x|\le\sigma_s\sigma_u$ define the Mehler
kernel and its order-$N$ truncation
\begin{equation}\label{eq:mehler_kernel}
\begin{split}
  K(s,u;x)&=\sum_{m\ge0}\hat c_m(s)\hat c_m(u)\Big(\frac{x}{\sigma_s\sigma_u}\Big)^m,\\
  K_N(s,u;x)&=\sum_{m=0}^N\hat c_m(s)\hat c_m(u)\Big(\frac{x}{\sigma_s\sigma_u}\Big)^m.
\end{split}
\end{equation}
At a point with $\sigma_s=0$, use $f(s,0)$ in order zero and zero tensors at positive orders.
Joint measurability of the scalar coefficients follows by Gaussian integration.
Each chaos tensor is strongly measurable in a separable tensor power, with norm bounded by
$\sup_s\|X_s\|_2/\sqrt{m!}$. The partial chaos sums are strongly measurable, and their pointwise
$L^2$ limit gives the stated Bochner integrand.

By Mehler's formula \citep{Mehler1866,Janson1997}, $K(s,u;x)=\E[f(s,U)f(u,V)]$ whenever $(U,V)$ is
centred bivariate Gaussian with variances $(\sigma_s^2,\sigma_u^2)$ and covariance $x$; in particular
$\E[X_sX_u]=K(s,u;R(s,u))$. For the Wick exponential $f(s,y)=\xi_0(s)\exp(\eta y-\tfrac12\eta^2\sigma_s^2)$,
which governs rough-Bergomi spot and forward variance, $\hat c_m(s)=\xi_0(s)(\eta\sigma_s)^m/\sqrt{m!}$ and
$K(s,u;x)=\xi_0(s)\xi_0(u)e^{\eta^2x}$. For mixed rough Bergomi, $f(s,y)=\xi_0(s)\sum_{r=1}^m\lambda_r\exp(\eta_r y-\tfrac12\eta_r^2\sigma_s^2)$
gives
\[
  K(s,u;x)=\xi_0(s)\xi_0(u)\sum_{r,q=1}^m\lambda_r\lambda_q e^{\eta_r\eta_q x}.
\]
For rough Stein--Stein integrated variance $f(s,y)=(\sigma_0+\nu y)^2$, $K$ is a degree-two polynomial in $x$.

\begin{definition}[Admissible Gaussian representation]\label{def:rep}
A map $s\mapsto g(s)\in\R^k$ with a standard Gaussian vector $\xi\in\R^k$ is an \emph{admissible
representation} of $Y$ if $g(s)\cdot\xi=W(\Pi e_s)$ for all $s\in\mathcal{I}$, where $\Pi$ is the
orthogonal projection of $\cH$ onto a closed subspace $S$ spanned by $\xi$. Its retained and residual
covariance kernels are
\[
  \RP(s,u)=g(s)\cdot g(u),\qquad D(s,u)=R(s,u)-\RP(s,u).
\]
\end{definition}

Because orthogonality implies independence for jointly Gaussian variables, Definition~\ref{def:rep} holds
if and only if $g(s)\cdot\xi=\E[Y_s\mid\xi]$ and the residual field $Y_s-g(s)\cdot\xi$ is independent of
$\xi$ with positive semidefinite covariance $D$ (the \hyperref[app:proofs]{proof appendix}). A non-orthogonal truncation
violates this contract and can produce indefinite $R-gg^\top$; every representation used below is
explicitly verified to be an orthogonal projection.

\subsection{Exact joint error identity and two-sided residual bounds}\label{sec:ledger}
Let $\Gamma(\Pi)$ denote the second quantisation of $\Pi$ and $J_{\le N}$ the orthogonal projection onto
Wiener chaoses of order at most $N\in\N_0\cup\{\infty\}$ \citep{Janson1997,Nualart2006}. They commute, so
$Q=J_{\le N}\Gamma(\Pi)$ is an orthogonal projection of $L^2(\Omega)$. Conditionally on $W(S)$,
$\Gamma(\Pi)X_s=\E[X_s\mid W(S)]=(S_{D(s,s)}f(s,\cdot))(W(\Pi e_s))$, where $S_v h(y)=\E[h(y+\sqrt{v}\,Z)]$
smooths $f$ by the residual variance $D(s,s)$. For a target measure $\mu$ and a monitoring measure $\nu$,
the combined surrogate is $Q\Lambda_\nu$.

\begin{theorem}[Exact joint error identity]\label{thm:master}
For finite signed Borel measures $\mu,\nu$ on $\mathcal{I}$, any admissible representation satisfying
Definition~\ref{def:rep}, and any $N\in\N_0\cup\{\infty\}$,
\begin{equation}\label{eq:master}
\begin{split}
  \norm{\Lambda_\mu-Q\Lambda_\nu}_2^2
  =\iint K(R)\dd\mu\dd\mu&-2\iint K_N(\RP)\dd\mu\dd\nu\\
  &+\iint K_N(\RP)\dd\nu\dd\nu,
\end{split}
\end{equation}
where $K(R)$ abbreviates $K(s,u;R(s,u))$ and $K_N(\RP)$ abbreviates $K_N(s,u;\RP(s,u))$.
\end{theorem}

Define the component terms
\begin{align*}
  P_\rho&=\norm{\Lambda_\rho-\Gamma(\Pi)\Lambda_\rho}_2^2\\
  &=\iint\big(K(R)-K(\RP)\big)\dd\rho\dd\rho,\\
  C_\rho&=\norm{\Gamma(\Pi)\Lambda_\rho-Q\Lambda_\rho}_2^2\\
  &=\iint\big(K(\RP)-K_N(\RP)\big)\dd\rho\dd\rho,\\
  M&=\iint K(R)\dd(\mu-\nu)\dd(\mu-\nu),\\
  M_Q&=\iint K_N(\RP)\dd(\mu-\nu)\dd(\mu-\nu),\\
  I&=2\iint\big(K(R)-K_N(\RP)\big)\dd(\mu-\nu)\dd\nu.
\end{align*}
These yield two equivalent decompositions of \eqref{eq:master}:

\begin{corollary}[Reconciled sequential and orthogonal ledgers]\label{cor:ledgers}
Under Theorem~\ref{thm:master},
\[
  \norm{\Lambda_\mu-Q\Lambda_\nu}_2^2=M+P_\nu+C_\nu+I=M_Q+P_\mu+C_\mu,
\]
where $M_Q,P_\mu,C_\mu\ge0$, $0\le M_Q\le M$, and
$|I|\le2\sqrt{(M-M_Q)(P_\nu+C_\nu)}$. In particular,
$P_\mu+C_\mu\le\norm{\Lambda_\mu-Q\Lambda_\nu}_2^2\le M+P_\mu+C_\mu$.
\end{corollary}

In the terminal case $\mu=\delta_T$ with $f(T,y)=A\exp^\diamond(\sigma y)\diamond(1+qy)$,
Theorem~\ref{thm:master} recovers the exact closed-form chaos truncation tail $R_N=\E|X-X_N|^2$ derived
in \citet{AHMED2026111045}; the present identity extends that exact error calculus to measure-indexed
integrated functionals and retains the cross-interaction $I$ between discretisation, projection, and
chaos truncation.

To diagnose in advance how a representation's residual covariance $D$ converts into output projection error
$P_\rho$, let $K'(s,u;x)=\partial_x K(s,u;x)$. By Price's theorem \citep{Price1958,Voigtlaender2017Price},
$K'(s,u;x)=\E[\partial_y f(s,U)\,\partial_y f(u,V)]$.

\begin{theorem}[Two-sided residual-to-output sandwich]\label{thm:sandwich}
Suppose $f(s,\cdot)$ is weakly differentiable with $\sup_{s\in\mathcal{I}}\norm{\partial_y f(s,Y_s)}_2<\infty$.
For every finite signed Borel measure $\rho$ and every admissible representation,
\begin{equation}\label{eq:sandwich}
\begin{split}
  \iint c_1(s)c_1(u)D(s,u)\dd\rho\dd\rho
  &\;\le\;\iint K'(\RP)D\dd\rho\dd\rho\\
  &\;\le\;P_\rho\;\le\;\iint K'(R)D\dd\rho\dd\rho,
\end{split}
\end{equation}
where $c_1(s)=\E[\partial_y f(s,Y_s)]$. The first two inequalities are equalities whenever $\Lambda_\rho$
lies in the sum of the zeroth and first Wiener chaoses.
\end{theorem}

The upper inequality is the subspace Gaussian Poincar\'e bound of \citet{ZahmConstantinePrieurMarzouk2020};
the lower inequalities follow from Schur-power monotonicity on the positive semidefinite interpolating path
$R_\theta=\RP+\theta D$ (the \hyperref[app:proofs]{proof appendix}) and extend to non-differentiable square-integrable
outputs as inequalities of integrated quadratic forms. In that case the middle lower bound means
the increasing limit of its finite-chaos forms and $c_1(s)=\hat c_1(s)/\sigma_s$ where $\sigma_s>0$ and zero otherwise.
At zero variance the associated residual vector is zero. This convention covers indicator outputs.

\subsection{Measure-axis convergence rate}\label{sec:rate}
When a continuous variance integral $\Lambda_\mu=\int_0^T f(s,Y_s)\,g(s)\dd s$ is approximated by an
$n$-cell Riemann sum $\Lambda_{\mu_n}$ with $\mu_n=\sum_{i=1}^n w_i\delta_{t_i}$ ($t_i\in\text{cell }i$,
$w_i=\int_{\text{cell}_i}g$), the exact monitoring error $M(\mu_n)=\norm{\Lambda_\mu-\Lambda_{\mu_n}}_2^2$
is governed by the local covariance regularity of $Y$.

\begin{assumption}[Covariance regularity]\label{ass:K}
On $[0,T]$, $R$ is continuous with $|R|\le\bar s$, and there exist constants $c_R,A_R,A_R'>0$ such that:
\emph{(K1)} $\norm{e_s-e_u}^2\le c_R|s-u|^{2H}$;
\emph{(K2)} for $0<s<u$ and $p=1/2-H$, the mixed derivative satisfies
\[
 |\partial_{su}R|\le A_R\{(u-s)^{-1-2p}+s^{-p}u^{-1-p}\}.
\]
The transposed bound holds on $0<u<s$. \emph{(K3)} on the open off-diagonal triangles,
$|\partial_sR|\le A'_R(s^{-2p}+|s-u|^{-2p})$.
The derivatives there are continuous. For $H=1/2$ we treat Brownian covariance separately.
\end{assumption}

For fractional Brownian motion the off-diagonal kernel derivatives follow directly from
$2R=s^{2H}+u^{2H}-|u-s|^{2H}$. The endpoint term in (K2) also covers the
Riemann--Liouville field, as shown in the \hyperref[app:proofs]{proof appendix}.

\begin{theorem}[Measure-axis distance rate]\label{thm:rate}
Let $f(s,y)=\xi_0\exp(\eta y-\tfrac12\eta^2\sigma_s^2)$, let
$g$ be bounded measurable on $[0,T]$, and let $\mu_n$ be any cell-matched discretisation on a uniform
partition of width $h=T/n$. Under Assumption~\ref{ass:K} with $0<H<1/2$, or for Brownian covariance at $H=1/2$,
\begin{equation}\label{eq:rate}
\begin{split}
  M(\mu_n)=\norm{\Lambda_\mu-\Lambda_{\mu_n}}_2^2
  \;\le\;&C(H,T,\norm{g}_\infty,\bar s,c_R,A_R,A_R')\\
  &\times\xi_0^2\,\eta^2 e^{\eta^2\bar s}(1+\eta^2\bar s)\,n^{-(1+2H)}.
\end{split}
\end{equation}
A finite linear combination has the same exponent by the triangle inequality in $L^2$, with the coefficient-dependent constant obtained by summing the individual RMS bounds. At the Brownian boundary $H=\tfrac12$ with $g\equiv1$ and left-point sampling,
$\lim_{n\to\infty} n^2 M(\mu_n)=\xi_0^2 T^2(e^{\eta^2 T}-1)/3$.
\end{theorem}

Theorem~\ref{thm:rate} (proved in the \hyperref[app:proofs]{proof appendix}) contrasts with the $O(n^{-1})$ variance
error observed for signed first-chaos energy discrepancies of deterministic integrands
\citep{AHMED2026111045}: because $M(\mu_n)$ is a squared $L^2$ distance rather
than a difference of squared norms, the rate $n^{-(1+2H)}$ holds uniformly across in-cell node placements
without endpoint vanishing conditions. The WIS solution approximation studied by \citet{ErdoganLordSchieven2024WIS} concerns a different error criterion. For $H=0.10$, the RMS monitoring error decays only as
$n^{-0.60}$, explaining quantitatively why daily grids leave substantial continuum bias in rough-volatility
variance swaps and why conditional expectation is essential.
